% GENERATED FILE. Edit canonical lessons, then run npm run generate:books.
% canonical-source-hash: fnv1a64:a1e21af784fd01fe
% canonical-lessons: HI-C353-samsad, HI-C353-sasan, HI-C353-niti, HI-C353-prastav, HI-C353-thahrana

\chapter{Who Governs}
\label{ch:hi-who-governs}
\hlchaptermodality{353}

\begin{chapteropening}
\textbf{By the end of this chapter:} \emph{I can say who makes the rules and who is to blame.}

Complete the last lesson of chapter 353: I can say who makes the rules and who is to blame.
\end{chapteropening}

\section[\texorpdfstring{\hi{संसद} (saṁsad)}{संसद (saṁsad)}]{\hi{संसद} (saṁsad) — parliament}

\label{lesson:HI-C353-samsad}



\begin{quote}
Before the new one: say the Hindi for guilty, then the Hindi for bail.
\end{quote}

\subsection*{You'll want to know: \hi{संसद}}
\textbf{\hi{संसद}} — \emph{saṁsad} — \textquotedblleft{}parliament\textquotedblright{}.

\begin{quote}
\hi{संसद} \hi{में} \hi{आज} \hi{नए} \hi{क़ानून} \hi{पर} \hi{बहस} \hi{हुई।}
\end{quote}

\begin{quote}
\emph{saṁsad meṁ āj nae qānūn par bahas huī.}
\end{quote}

\begin{quote}
Parliament debated the new law today.
\end{quote}

Feminine.

\begin{cousinweb}
\textbf{\hi{संसद}} (\emph{saṁsad}) is a tatsama, borrowed from Sanskrit \emph{saṃsad}, \textquotedblleft{}an assembly\textquotedblright{} (Monier-Williams).
\end{cousinweb}

\begin{grammarlens}[title={What you've built}]
One more word for giving reasons and opinions.
\end{grammarlens}

\subsection*{Your turn}
\begin{itemize}
\raggedright
  \item \emph{Say it:} \emph{saṁsad}
  \item \emph{Say it:} \emph{saṁsad}, once more
  \item \emph{Say it:} \emph{saṁsad meṁ āj nae qānūn par bahas huī.} — Parliament debated the new law today.
  \item \emph{Recall:} say the Hindi for guilty, then the Hindi for bail, then say \emph{saṁsad} again
\end{itemize}

\begin{tcolorbox}[breakable,colback=teal!4,colframe=teal!35!black,title={Before you move on}]
What does \hi{संसद} mean? (\textquotedblleft{}Parliament\textquotedblright{}.) Say it once more.
\end{tcolorbox}
\section[\texorpdfstring{\hi{शासन} (śāsan)}{शासन (śāsan)}]{\hi{शासन} (śāsan) — rule, governance}

\label{lesson:HI-C353-sasan}



\begin{quote}
Before the new one: say the Hindi for bail, then the Hindi for parliament.
\end{quote}

\subsection*{You'll want to know: \hi{शासन}}
\textbf{\hi{शासन}} — \emph{śāsan} — \textquotedblleft{}rule, governance\textquotedblright{}.

\begin{quote}
\hi{राजा} \hi{का} \hi{शासन} \hi{चालीस} \hi{साल} \hi{चला।}
\end{quote}

\begin{quote}
\emph{rājā kā śāsan cālīs sāl calā.}
\end{quote}

\begin{quote}
The king's rule lasted forty years.
\end{quote}

Masculine.

\begin{cousinweb}
\textbf{\hi{शासन}} (\emph{śāsan}) is a tatsama, borrowed from Sanskrit \emph{śāsana}, \textquotedblleft{}governing; instruction\textquotedblright{} (Monier-Williams).
\end{cousinweb}

\begin{grammarlens}[title={What you've built}]
One more word for giving reasons and opinions.
\end{grammarlens}

\subsection*{Your turn}
\begin{itemize}
\raggedright
  \item \emph{Say it:} \emph{śāsan}
  \item \emph{Say it:} \emph{śāsan}, once more
  \item \emph{Say it:} \emph{rājā kā śāsan cālīs sāl calā.} — The king's rule lasted forty years.
  \item \emph{Recall:} say the Hindi for bail, then the Hindi for parliament, then say \emph{śāsan} again
\end{itemize}

\begin{tcolorbox}[breakable,colback=teal!4,colframe=teal!35!black,title={Before you move on}]
What does \hi{शासन} mean? (\textquotedblleft{}Rule, governance\textquotedblright{}.) Say it once more.
\end{tcolorbox}
\section[\texorpdfstring{\hi{नीति} (nīti)}{नीति (nīti)}]{\hi{नीति} (nīti) — a policy}

\label{lesson:HI-C353-niti}



\begin{quote}
Before the new one: say the Hindi for parliament, then the Hindi for governance.
\end{quote}

\subsection*{You'll want to know: \hi{नीति}}
\textbf{\hi{नीति}} — \emph{nīti} — \textquotedblleft{}a policy\textquotedblright{}.

\begin{quote}
\hi{सरकार} \hi{की} \hi{नई} \hi{नीति} \hi{किसानों} \hi{के} \hi{लिए} \hi{अच्छी} \hi{है।}
\end{quote}

\begin{quote}
\emph{sarkār kī naī nīti kisānoṁ ke liye acchī hai.}
\end{quote}

\begin{quote}
The government's new policy is good for farmers.
\end{quote}

Feminine.

\begin{cousinweb}
\textbf{\hi{नीति}} (\emph{nīti}) is a tatsama, borrowed from Sanskrit \emph{nīti}, \textquotedblleft{}conduct, policy\textquotedblright{} (Monier-Williams).
\end{cousinweb}

\begin{grammarlens}[title={What you've built}]
One more word for giving reasons and opinions.
\end{grammarlens}

\subsection*{Your turn}
\begin{itemize}
\raggedright
  \item \emph{Say it:} \emph{nīti}
  \item \emph{Say it:} \emph{nīti}, once more
  \item \emph{Say it:} \emph{sarkār kī naī nīti kisānoṁ ke liye acchī hai.} — The government's new policy is good for farmers.
  \item \emph{Recall:} say the Hindi for parliament, then the Hindi for governance, then say \emph{nīti} again
\end{itemize}

\begin{tcolorbox}[breakable,colback=teal!4,colframe=teal!35!black,title={Before you move on}]
What does \hi{नीति} mean? (\textquotedblleft{}A policy\textquotedblright{}.) Say it once more.
\end{tcolorbox}
\section[\texorpdfstring{\hi{प्रस्ताव} (prastāv)}{प्रस्ताव (prastāv)}]{\hi{प्रस्ताव} (prastāv) — a proposal, a motion}

\label{lesson:HI-C353-prastav}



\begin{quote}
Before the new one: say the Hindi for governance, then the Hindi for a policy.
\end{quote}

\subsection*{You'll want to know: \hi{प्रस्ताव}}
\textbf{\hi{प्रस्ताव}} — \emph{prastāv} — \textquotedblleft{}a proposal, a motion\textquotedblright{}.

\begin{quote}
\hi{मंत्री} \hi{ने} \hi{संसद} \hi{में} \hi{नया} \hi{प्रस्ताव} \hi{रखा।}
\end{quote}

\begin{quote}
\emph{mantrī ne saṁsad meṁ nayā prastāv rakhā.}
\end{quote}

\begin{quote}
The minister put a new motion before parliament.
\end{quote}

Masculine.

\begin{cousinweb}
\textbf{\hi{प्रस्ताव}} (\emph{prastāv}) is a tatsama, borrowed from Sanskrit \emph{prastāva}, \textquotedblleft{}an introduction; the subject of a talk\textquotedblright{} (Monier-Williams). Monier-Williams does not give the sense of a formal proposal, which is usually taken as a modern Hindi one.
\end{cousinweb}

\begin{grammarlens}[title={What you've built}]
One more word for giving reasons and opinions.
\end{grammarlens}

\subsection*{Your turn}
\begin{itemize}
\raggedright
  \item \emph{Say it:} \emph{prastāv}
  \item \emph{Say it:} \emph{prastāv}, once more
  \item \emph{Say it:} \emph{mantrī ne saṁsad meṁ nayā prastāv rakhā.} — The minister put a new motion before parliament.
  \item \emph{Recall:} say the Hindi for governance, then the Hindi for a policy, then say \emph{prastāv} again
\end{itemize}

\begin{tcolorbox}[breakable,colback=teal!4,colframe=teal!35!black,title={Before you move on}]
What does \hi{प्रस्ताव} mean? (\textquotedblleft{}A proposal, a motion\textquotedblright{}.) Say it once more.
\end{tcolorbox}
\section[\texorpdfstring{\hi{ठहराना} (ṭhahrānā)}{ठहराना (ṭhahrānā)}]{\hi{ठहराना} (ṭhahrānā) — to hold to be, to declare (guilty, responsible)}

\label{lesson:HI-C353-thahrana}



\begin{quote}
Before the new one: say the Hindi for a policy, then the Hindi for a proposal.
\end{quote}

\subsection*{You'll want to know: \hi{ठहराना}}
\textbf{\hi{ठहराना}} — \emph{ṭhahrānā} — \textquotedblleft{}to hold to be, to declare (guilty, responsible)\textquotedblright{}.

\begin{quote}
\hi{लोगों} \hi{ने} \hi{सरकार} \hi{को} \hi{ज़िम्मेदार} \hi{ठहराया।}
\end{quote}

\begin{quote}
\emph{logoṁ ne sarkār ko zimmedār ṭhahrāyā.}
\end{quote}

\begin{quote}
People held the government responsible.
\end{quote}

Past, with \hi{ने} (he or she …): \hi{उसने} \hi{मुझे} \hi{दोषी} \hi{ठहराया}, \emph{usne mujhe doṣī ṭhahrāyā}.

\begin{cousinweb}
\textbf{\hi{ठहराना}} (\emph{ṭhahrānā}) is the causative of \hi{ठहरना}, \textquotedblleft{}to stop, to settle\textquotedblright{}, which CDIAL traces to a reconstructed \emph{*stabhira} (CDIAL 13680): to make a judgement settle on someone.
\end{cousinweb}

\begin{grammarlens}[title={What you've built}]
One more word for giving reasons and opinions.
\end{grammarlens}

\subsection*{Your turn}
\begin{itemize}
\raggedright
  \item \emph{Say it:} \emph{ṭhahrānā}
  \item \emph{Say it:} \emph{ṭhahrānā}, once more
  \item \emph{Say it:} \emph{logoṁ ne sarkār ko zimmedār ṭhahrāyā.} — People held the government responsible.
  \item \emph{Recall:} say the Hindi for a policy, then the Hindi for a proposal, then say \emph{ṭhahrānā} again
\end{itemize}

\begin{tcolorbox}[breakable,colback=teal!4,colframe=teal!35!black,title={Before you move on}]
What does \hi{ठहराना} mean? (\textquotedblleft{}To hold to be, to declare (guilty, responsible)\textquotedblright{}.) Say it once more.
\end{tcolorbox}
